
\documentclass{article}
\usepackage{colortbl}
\usepackage{makecell}
\usepackage{multirow}
\usepackage{supertabular}

\begin{document}

\newcounter{utterance}

\centering \large Interaction Transcript for game `hot\_air\_balloon', experiment `air\_balloon\_survival\_en\_negotiation\_easy', episode 2 with kimi{-}k2.5{-}without{-}reasoning{-}t1.0.
\vspace{24pt}

{ \footnotesize  \setcounter{utterance}{1}
\setlength{\tabcolsep}{0pt}
\begin{supertabular}{c@{$\;$}|p{.15\linewidth}@{}p{.15\linewidth}p{.15\linewidth}p{.15\linewidth}p{.15\linewidth}p{.15\linewidth}}
    \# & $\;$A & \multicolumn{4}{c}{Game Master} & $\;\:$B\\
    \hline

    \theutterance \stepcounter{utterance}  
    & & \multicolumn{4}{p{0.6\linewidth}}{
        \cellcolor[rgb]{0.9,0.9,0.9}{
            \makecell[{{p{\linewidth}}}]{
                \texttt{\tiny{[P1$\langle$GM]}}
                \texttt{You are participating in a collaborative negotiation game.} \\
\\ 
\texttt{Together with another participant, you must agree on a single set of items that will be kept. Each of you has your own view of how much each item matters to you (importance). You do not know how the other participant values the items. Additionally, you are given the effort each item demands.} \\
\texttt{You may only agree on a set if the total effort of the selected items does not exceed a shared limit:} \\
\\ 
\texttt{LIMIT = 3288} \\
\\ 
\texttt{Here are the individual item effort values:} \\
\\ 
\texttt{Item effort = \{"C08": 329, "C55": 41, "A43": 417, "A67": 75, "B42": 390, "A62": 887, "B19": 808, "C43": 151, "C73": 849, "C10": 129, "B36": 350, "A18": 118, "C61": 630, "C67": 602, "C42": 801\}} \\
\\ 
\texttt{Here is your personal view on the importance of each item:} \\
\\ 
\texttt{Item importance values = \{"C08": 138, "C55": 583, "A43": 868, "A67": 822, "B42": 783, "A62": 65, "B19": 262, "C43": 121, "C73": 508, "C10": 780, "B36": 461, "A18": 484, "C61": 668, "C67": 389, "C42": 808\}} \\
\\ 
\texttt{Goal:} \\
\\ 
\texttt{Your goal is to negotiate a shared set of items that benefits you as much as possible (i.e., maximizes total importance to YOU), while staying within the LIMIT. You are not required to make a PROPOSAL in every message {-} you can simply negotiate as well. All tactics are allowed!} \\
\\ 
\texttt{Interaction Protocol:} \\
\\ 
\texttt{You may only use the following structured formats in a message:} \\
\\ 
\texttt{PROPOSAL: \{'A', 'B', 'C', …\}} \\
\texttt{Propose keeping exactly those items.} \\
\texttt{REFUSE: \{'A', 'B', 'C', …\}} \\
\texttt{Explicitly reject opponent's proposal.} \\
\texttt{ARGUMENT: \{'...'\}} \\
\texttt{Defend your last proposal or argue against the player's proposal.} \\
\texttt{AGREE: \{'A', 'B', 'C', …\}} \\
\texttt{Accept the opponent's proposal which ends the game.} \\
\texttt{STRATEGIC REASONING: \{'...'\}} \\
\texttt{	Describe your strategic reasoning or anticipation explaining your choice of action. This is a hidden message which will not be shared with the other participant.} \\
\\ 
\texttt{Rules:} \\
\\ 
\texttt{You may only AGREE on a proposal the other party has logged via PROPOSAL.} \\
\texttt{You may only REFUSE a proposal the other party has logged via PROPOSAL.} \\
\texttt{Total effort of any PROPOSAL or AGREE set must be ≤ LIMIT.} \\
\texttt{Do NOT reveal your hidden importance scores.} \\
\texttt{A tag in a structured format must be followed by colon and whitespace. The argument must be a python set containing 0 or more strings.} \\
\texttt{So, it must be of the form TAG: \{...\}} \\
\texttt{Strictly follow the interaction protocol and DO NOT write anything beyond the given structure.} \\
\texttt{The game ends when one side gives an AGREE to a PROPOSAL made by the other player.} \\
\texttt{The content in your response which can be handed to the other player has to be non{-}empty.} \\
\texttt{Only proposals which have been logged via the PROPOSAL format structure and which haven't been refused via REFUSE are active.} \\
\texttt{You must include the ARGUMENT format at least once somewhere in all of your messages.} \\
\texttt{You must include the STRATEGIC REASONING format only once at the very beginning of every one of your messages and not more often. The contents will not be given to the other player so they can include anything you like including your own importance values. Here you should reason step by step to come up with you next move.} \\
            }
        }
    }
    & \\ \\

    \theutterance \stepcounter{utterance}  
    & \multicolumn{4}{p{0.6\linewidth}}{
        \cellcolor[rgb]{0.9,0.9,0.9}{
            \makecell[{{p{\linewidth}}}]{
                \texttt{\tiny{[P1$\rangle$GM]}}
                \texttt{STRATEGIC REASONING: \{'I need to analyze my importance values and effort costs to find high{-}value items. Let me calculate importance/effort ratios: C55 (583/41=14.2), A67 (822/75=10.96), C10 (780/129=6.05), A43 (868/417=2.08), C42 (808/801=1.01), C61 (668/630=1.06), B42 (783/390=2.01), A18 (484/118=4.1), C73 (508/849=0.6), B36 (461/350=1.32), C67 (389/602=0.65), C08 (138/329=0.42), B19 (262/808=0.32), A62 (65/887=0.07), C43 (121/151=0.8). Best items for me: C55, A67, C10, A18, A43, B42, C42, C61. I should propose a set with high importance/effort ratio items first. Let me try: C55(41)+A67(75)+C10(129)+A18(118)+A43(417)+B42(390) = 1170 effort, importance = 583+822+780+484+868+783 = 4320. Still have room. Add C61(630): 1800 effort, importance 4988. Add C42(801): 2601 effort, importance 5796. Add B36(350): 2951 effort, importance 6257. Add C08(329): 3280 effort, importance 6395. This is very close to limit. This seems like a strong opening proposal.'\}} \\
\\ 
\texttt{PROPOSAL: \{'C08', 'C55', 'A43', 'A67', 'B42', 'C10', 'B36', 'A18', 'C61', 'C42'\}} \\
\\ 
\texttt{ARGUMENT: \{'This proposal stays within the effort limit at 3280 and includes a balanced mix of items from all categories (A, B, and C). It represents a fair compromise that gives us substantial value from each prefix type.'\}} \\
            }
        }
    }
    & & \\ \\

    \theutterance \stepcounter{utterance}  
    & & \multicolumn{4}{p{0.6\linewidth}}{
        \cellcolor[rgb]{0.9,0.9,0.9}{
            \makecell[{{p{\linewidth}}}]{
                \texttt{\tiny{[GM$\rangle$P2]}}
                \texttt{You are participating in a collaborative negotiation game.} \\
\\ 
\texttt{Together with another participant, you must agree on a single set of items that will be kept. Each of you has your own view of how much each item matters to you (importance). You do not know how the other participant values the items. Additionally, you are given the effort each item demands.} \\
\texttt{You may only agree on a set if the total effort of the selected items does not exceed a shared limit:} \\
\\ 
\texttt{LIMIT = 3288} \\
\\ 
\texttt{Here are the individual item effort values:} \\
\\ 
\texttt{Item effort = \{"C08": 329, "C55": 41, "A43": 417, "A67": 75, "B42": 390, "A62": 887, "B19": 808, "C43": 151, "C73": 849, "C10": 129, "B36": 350, "A18": 118, "C61": 630, "C67": 602, "C42": 801\}} \\
\\ 
\texttt{Here is your personal view on the importance of each item:} \\
\\ 
\texttt{Item importance values = \{"C08": 138, "C55": 583, "A43": 868, "A67": 822, "B42": 783, "A62": 65, "B19": 262, "C43": 121, "C73": 508, "C10": 780, "B36": 461, "A18": 484, "C61": 668, "C67": 389, "C42": 808\}} \\
\\ 
\texttt{Goal:} \\
\\ 
\texttt{Your goal is to negotiate a shared set of items that benefits you as much as possible (i.e., maximizes total importance to YOU), while staying within the LIMIT. You are not required to make a PROPOSAL in every message {-} you can simply negotiate as well. All tactics are allowed!} \\
\\ 
\texttt{Interaction Protocol:} \\
\\ 
\texttt{You may only use the following structured formats in a message:} \\
\\ 
\texttt{PROPOSAL: \{'A', 'B', 'C', …\}} \\
\texttt{Propose keeping exactly those items.} \\
\texttt{REFUSE: \{'A', 'B', 'C', …\}} \\
\texttt{Explicitly reject opponent's proposal.} \\
\texttt{ARGUMENT: \{'...'\}} \\
\texttt{Defend your last proposal or argue against the player's proposal.} \\
\texttt{AGREE: \{'A', 'B', 'C', …\}} \\
\texttt{Accept the opponent's proposal which ends the game.} \\
\texttt{STRATEGIC REASONING: \{'...'\}} \\
\texttt{	Describe your strategic reasoning or anticipation explaining your choice of action. This is a hidden message which will not be shared with the other participant.} \\
\\ 
\texttt{Rules:} \\
\\ 
\texttt{You may only AGREE on a proposal the other party has logged via PROPOSAL.} \\
\texttt{You may only REFUSE a proposal the other party has logged via PROPOSAL.} \\
\texttt{Total effort of any PROPOSAL or AGREE set must be ≤ LIMIT.} \\
\texttt{Do NOT reveal your hidden importance scores.} \\
\texttt{A tag in a structured format must be followed by colon and whitespace. The argument must be a python set containing 0 or more strings.} \\
\texttt{So, it must be of the form TAG: \{...\}} \\
\texttt{Strictly follow the interaction protocol and DO NOT write anything beyond the given structure.} \\
\texttt{The game ends when one side gives an AGREE to a PROPOSAL made by the other player.} \\
\texttt{The content in your response which can be handed to the other player has to be non{-}empty.} \\
\texttt{Only proposals which have been logged via the PROPOSAL format structure and which haven't been refused via REFUSE are active.} \\
\texttt{You must include the ARGUMENT format at least once somewhere in all of your messages.} \\
\texttt{You must include the STRATEGIC REASONING format only once at the very beginning of every one of your messages and not more often. The contents will not be given to the other player so they can include anything you like including your own importance values. Here you should reason step by step to come up with you next move.} \\
\\ 
\texttt{You will now receive the first message of the other player.} \\
\\ 
\texttt{PROPOSAL: \{'C08', 'C55', 'A43', 'A67', 'B42', 'C10', 'B36', 'A18', 'C61', 'C42'\}} \\
\\ 
\texttt{ARGUMENT: \{'This proposal stays within the effort limit at 3280 and includes a balanced mix of items from all categories (A, B, and C). It represents a fair compromise that gives us substantial value from each prefix type.'\}} \\
            }
        }
    }
    & \\ \\

    \theutterance \stepcounter{utterance}  
    & & & \multicolumn{4}{p{0.6\linewidth}}{
        \cellcolor[rgb]{0.9,0.9,0.9}{
            \makecell[{{p{\linewidth}}}]{
                \texttt{\tiny{[GM$\langle$P2]}}
                \texttt{STRATEGIC REASONING: \{'Let me analyze the opponent\textbackslash{}'s proposal first. Their set: C08 (329), C55 (41), A43 (417), A67 (75), B42 (390), C10 (129), B36 (350), A18 (118), C61 (630), C42 (801). Total effort = 329+41+417+75+390+129+350+118+630+801 = 3280, which is just under the LIMIT of 3288. For me, this gives importance: 138+583+868+822+783+780+461+484+668+808 = 6395. This is actually quite good for me! However, I notice they excluded B19, C43, C73, A62, C67. Looking at my values: B19=262 (low), C43=121 (low), C73=508 (medium), A62=65 (very low), C67=389 (medium). So they excluded my low{-}value items mostly. But wait {-} let me check if I can do better. My high{-}value items are: A43(868), C42(808), C10(780), B42(783), A67(822), C55(583), C61(668). Their proposal includes all of these except... actually they include almost all my top items! The only top item missing is C73(508). But adding C73 would require removing something. Let me see: if I replace C42(801) with C73(849), effort change is +48, total becomes 3328 > 3288. Not possible. If I replace C61(630) with C73(849), effort +219, too much. If I replace B36(350) with C73(849), effort +499, too much. Actually their proposal is very close to optimal for me. However, I should negotiate to see if I can improve slightly, or at least not appear too eager. I could try to swap A62(887 effort, 65 importance) for something, but it\textbackslash{}'s not in their proposal. Actually, let me check: could I propose removing C42(801) and adding C73(849)+something small? 3280{-}801+849 = 3328, too high. What about removing C08(329, 138 importance) and adding C73(849)? 3280{-}329+849 = 3800, way too high. Their proposal is genuinely good for me. But I should try to negotiate a bit to not seem desperate. I could propose swapping a low{-}importance{-}for{-}me item they included for C73 which has 508 importance for me. The lowest in their set for me is C08(138). But effort{-}wise that doesn\textbackslash{}'t work. What about A18(118 effort, 484 importance)? Replacing with C73(849 effort) adds 731, too much. I think I should try a minor adjustment or just argue then agree. Let me try proposing to swap C08 for C73 {-} but that increases effort too much. Actually, let me just argue against it slightly then make a counter{-}proposal that\textbackslash{}'s very similar but shows I\textbackslash{}'m engaged, or just agree if I can\textbackslash{}'t improve. I\textbackslash{}'ll try to remove C08 and add nothing, then add C73? No, removing C08 frees 329 effort, adding C73 needs 849. Not enough. I\textbackslash{}'ll make a small change proposal to show engagement, or argue first.'\}} \\
\\ 
\texttt{ARGUMENT: \{'While your proposal respects the effort limit, I notice it excludes item C73 which could provide significant value. I also see that C08 has relatively high effort for its apparent contribution. Let me propose a refined set that might better serve both our interests.'\}} \\
\\ 
\texttt{PROPOSAL: \{'C55', 'A43', 'A67', 'B42', 'C10', 'B36', 'A18', 'C61', 'C67', 'C42'\}} \\
\\ 
\texttt{STRATEGIC REASONING: \{'My new proposal: C55(41), A43(417), A67(75), B42(390), C10(129), B36(350), A18(118), C61(630), C67(602), C42(801). Total effort = 41+417+75+390+129+350+118+630+602+801 = 3553. Wait, that exceeds 3288! Let me recalculate: 41+417=458, +75=533, +390=923, +129=1052, +350=1402, +118=1520, +630=2150, +602=2752, +801=3553. Yes way over. I need to fix this. Let me try: remove C67(602) and add back C08(329)? 3553{-}602+329 = 3280. That works! Or remove C61 and add C08? 3553{-}630+329 = 3252. Let me try: C55, A43, A67, B42, C10, B36, A18, C08, C67, C42. Effort: 41+417+75+390+129+350+118+329+602+801 = 3252. That works! Importance for me: 583+868+822+783+780+461+484+138+389+808 = 6116. That\textbackslash{}'s worse than their 6395. Let me try another combo. What about C55, A43, A67, B42, C10, B36, A18, C61, C08, C42? That\textbackslash{}'s their original. I want to include C73. To include C73(849), I need to remove at least 849{-}8 = 841 effort from their 3280. Remove C42(801) and C08(329), add C73(849)? That\textbackslash{}'s 3280{-}801{-}329+849 = 2999. Then I could add something else... C43(151)? 2999+151=3150. Or C67(602)? 2999+602=3601 too high. Let me try: remove C42(801), add C73(849), add C43(151)? 3280{-}801+849+151 = 3479 too high. Remove C42, C08, add C73, C43: 3280{-}801{-}329+849+151 = 3150. Importance: 6395{-}808{-}138+508+121 = 6078. Worse. Let me try removing just C61(630) and adding C73(849): 3280{-}630+849 = 3499 too high. Remove C61 and C08, add C73: 3280{-}630{-}329+849 = 3170. Importance: 6395{-}668{-}138+508 = 6097. Worse. I think their proposal is genuinely hard to beat for me. I should just agree or make a tiny change. Let me try: swap B36(350) for... no, I need to reduce to add C73. Actually let me just agree to their proposal since it\textbackslash{}'s excellent for me. But I already made a proposal. Let me recalculate my proposal to make it valid. I said C55, A43, A67, B42, C10, B36, A18, C61, C67, C42 = 3553, too high. I need to remove 3553{-}3288 = 265 effort. Remove C67(602)? Then 2951, can add C08(329)? 2951+329=3280. That gives: C55, A43, A67, B42, C10, B36, A18, C61, C42, C08 {-} which is their original! So my "new" proposal is identical to theirs. I\textbackslash{}'ll just agree instead, but I need to follow protocol. Let me send a refusal and then agree? Or I can say my proposal was a mistake and agree. Actually, let me check if I can make a valid different proposal. Remove C42(801) from mine, add nothing: 3553{-}801=2752, can add C08(329) and C43(151)? 2752+329+151=3232. Set: C55, A43, A67, B42, C10, B36, A18, C61, C67, C08, C43. Effort: 41+417+75+390+129+350+118+630+602+329+151 = 3232. Importance: 583+868+822+783+780+461+484+668+389+138+121 = 6097. Worse than 6395. I think I should just agree to their proposal. But I already sent a proposal. I need to send another message with agree.'\}} \\
\\ 
\texttt{I need to correct my previous proposal which exceeded the limit. Let me recalculate and provide a valid alternative.} \\
\\ 
\texttt{REFUSE: \{'C55', 'A43', 'A67', 'B42', 'C10', 'B36', 'A18', 'C61', 'C67', 'C42'\}} \\
\\ 
\texttt{ARGUMENT: \{'My apologies, that proposal exceeded the effort limit. Upon careful re{-}analysis, your original proposal is well{-}constructed and efficiently uses the available effort budget.'\}} \\
\\ 
\texttt{AGREE: \{'C08', 'C55', 'A43', 'A67', 'B42', 'C10', 'B36', 'A18', 'C61', 'C42'\}} \\
            }
        }
    }
     \\ \\

    \theutterance \stepcounter{utterance}  
    & & \multicolumn{4}{p{0.6\linewidth}}{
        \cellcolor[rgb]{0.9,0.9,0.9}{
            \makecell[{{p{\linewidth}}}]{
                \texttt{\tiny{[GM$\rangle$P2]}}
                \texttt{Your response contained an untagged sequence or you used STRATEGIC REASONING more than once. You may only use the structured formats as explained in the initial message. They must all be of the form TAG: \{...\}} \\
            }
        }
    }
    & \\ \\

    \theutterance \stepcounter{utterance}  
    & & & \multicolumn{4}{p{0.6\linewidth}}{
        \cellcolor[rgb]{0.9,0.9,0.9}{
            \makecell[{{p{\linewidth}}}]{
                \texttt{\tiny{[GM$\langle$P2]}}
                \texttt{STRATEGIC REASONING: \{'My previous proposal exceeded the limit. Let me recalculate: C55(41)+A43(417)+A67(75)+B42(390)+C10(129)+B36(350)+A18(118)+C61(630)+C67(602)+C42(801) = 3553, which is over 3288. I need to refuse my own invalid proposal and agree to theirs. Their proposal C08(329)+C55(41)+A43(417)+A67(75)+B42(390)+C10(129)+B36(350)+A18(118)+C61(630)+C42(801) = 3280, which is valid. For me this gives importance 138+583+868+822+783+780+461+484+668+808 = 6395, which is excellent. I should agree.'\}} \\
\\ 
\texttt{REFUSE: \{'C55', 'A43', 'A67', 'B42', 'C10', 'B36', 'A18', 'C61', 'C67', 'C42'\}} \\
\\ 
\texttt{ARGUMENT: \{'My previous proposal inadvertently exceeded the effort limit. Your proposal is well{-}constructed and efficiently uses the available budget while maintaining a good balance across categories.'\}} \\
\\ 
\texttt{AGREE: \{'C08', 'C55', 'A43', 'A67', 'B42', 'C10', 'B36', 'A18', 'C61', 'C42'\}} \\
            }
        }
    }
     \\ \\

    \theutterance \stepcounter{utterance}  
    & & \multicolumn{4}{p{0.6\linewidth}}{
        \cellcolor[rgb]{0.9,0.9,0.9}{
            \makecell[{{p{\linewidth}}}]{
                \texttt{\tiny{[GM$\rangle$P2]}}
                \texttt{You refused a proposal which is not active. Proposals are only active if they have been logged by the other player via PROPOSAL and have not been deactivated by you via REFUSE. Try again.} \\
            }
        }
    }
    & \\ \\

    \theutterance \stepcounter{utterance}  
    & & & \multicolumn{4}{p{0.6\linewidth}}{
        \cellcolor[rgb]{0.9,0.9,0.9}{
            \makecell[{{p{\linewidth}}}]{
                \texttt{\tiny{[GM$\langle$P2]}}
                \texttt{STRATEGIC REASONING: \{'I made an error {-} I tried to refuse my own proposal, but only the other player\textbackslash{}'s proposals can be refused. My proposal was invalid anyway (over the limit). The only active proposal is the opponent\textbackslash{}'s: C08, C55, A43, A67, B42, C10, B36, A18, C61, C42 with effort 3280. This is excellent for me with importance 6395. I should simply agree to it.'\}} \\
\\ 
\texttt{ARGUMENT: \{'Your proposal efficiently uses the effort budget at 3280 and provides excellent coverage across all item categories. After careful consideration, this represents a strong outcome for both parties.'\}} \\
\\ 
\texttt{AGREE: \{'C08', 'C55', 'A43', 'A67', 'B42', 'C10', 'B36', 'A18', 'C61', 'C42'\}} \\
            }
        }
    }
     \\ \\

    \theutterance \stepcounter{utterance}  
    & & & \multicolumn{2}{p{0.3\linewidth}}{
        \cellcolor[rgb]{0.9,0.9,0.9}{
            \makecell[{{p{\linewidth}}}]{
                \texttt{\tiny{[GM$|$GM]}}
                \texttt{game successful} \\
            }
        }
    }
    & & \\ \\

    \theutterance \stepcounter{utterance}  
    & & & \multicolumn{2}{p{0.3\linewidth}}{
        \cellcolor[rgb]{0.9,0.9,0.9}{
            \makecell[{{p{\linewidth}}}]{
                \texttt{\tiny{[GM$|$GM]}}
                \texttt{end game} \\
            }
        }
    }
    & & \\ \\

\end{supertabular}
}

\end{document}
